%%%
% Plantilla de Libro
% Modificación de una plantilla de Latex de Mathias Legrand (legrand.mathias@gmail.com)
% sobre modificaciones de Vel (vel@latextemplates.com) para adaptarla 
% al castellano y a las necesidades de escribir informática y matemáticas.
%
% Modificada por: Mario Román
%
% License:
% CC BY-NC-SA 3.0 (http://creativecommons.org/licenses/by-nc-sa/3.0/)
%%%

%%%%%%%%%%%%%%%%%%%%%%%%%%%%%%%%%%%%%%%%%
% The Legrand Orange Book
% LaTeX Template
% Version 2.0 (9/2/15)
%
% This template has been downloaded from:
% http://www.LaTeXTemplates.com
%
% Mathias Legrand (legrand.mathias@gmail.com) with modifications by:
% Vel (vel@latextemplates.com)
%
% License:
% CC BY-NC-SA 3.0 (http://creativecommons.org/licenses/by-nc-sa/3.0/)
%
% Compiling this template:
% This template uses biber for its bibliography and makeindex for its index.
% When you first open the template, compile it from the command line with the 
% commands below to make sure your LaTeX distribution is configured correctly:
%
% 1) pdflatex main
% 2) makeindex main.idx -s StyleInd.ist
% 3) biber main
% 4) pdflatex main x 2
%
% After this, when you wish to update the bibliography/index use the appropriate
% command above and make sure to compile with pdflatex several times 
% afterwards to propagate your changes to the document.
%
% This template also uses a number of packages which may need to be
% updated to the newest versions for the template to compile. It is strongly
% recommended you update your LaTeX distribution if you have any
% compilation errors.
%
% Important note:
% Chapter heading images should have a 2:1 width:height ratio,
% e.g. 920px width and 460px height.
%
%%%%%%%%%%%%%%%%%%%%%%%%%%%%%%%%%%%%%%%%%

%----------------------------------------------------------------------------------------
%	PACKAGES AND OTHER DOCUMENT CONFIGURATIONS
%----------------------------------------------------------------------------------------

%%% Configuración del papel.
% fleqn: Alinea ecuaciones a la izquierda
\documentclass[11pt, fleqn, spanish]{book}

%%% Castellano.
% noquoting: Permite uso de comillas no españolas.
% lcroman: Permite la enumeración con numerales romanos en minúscula.
% fontenc: Usa la fuente completa para que pueda copiarse correctamente del pdf.
\usepackage[spanish,es-noquoting,es-lcroman]{babel}
\usepackage[utf8]{inputenc}
\usepackage[T1]{fontenc}
\selectlanguage{spanish}

%%% Matemáticas.
% Paquetes de la AMS. Para entornos de ecuaciones.
\usepackage{amsmath,amsfonts,amsthm}

% Paquete propio para escribir teoría de categorías
\usepackage{categoriesdiag}


% Conmutative diagrams
%\usepackage{tkz-graph}
%\usetikzlibrary{arrows}
\usepackage{tikz}
\usetikzlibrary{matrix,arrows}
\tikzset{every loop/.style={min distance=10mm,in=0,out=60,looseness=10}}
\usepackage{tikz-cd}
\usetikzlibrary{cd}

%% Multicolumns
\usepackage{multicol}

%%% Shortcuts
\newcommand{\C}{\mathcal{C} }

%%% Twopart definition
\newcommand{\twopartdef}[4]
{
	\left\{
		\begin{array}{ll}
			#1 & \mbox{si } #2 \\
			#3 & \mbox{si } #4
		\end{array}
	\right.
}


%----------------------------------------------------------------------------------------

%%%%%%%%%%%%%%%%%%%%%%%%%%%%%%%%%%%%%%%%%
% The Legrand Orange Book
% Structural Definitions File
% Version 2.0 (9/2/15)
%
% Original author:
% Mathias Legrand (legrand.mathias@gmail.com) with modifications by:
% Vel (vel@latextemplates.com)
% 
% This file has been downloaded from:
% http://www.LaTeXTemplates.com
%
% License:
% CC BY-NC-SA 3.0 (http://creativecommons.org/licenses/by-nc-sa/3.0/)
%
%%%%%%%%%%%%%%%%%%%%%%%%%%%%%%%%%%%%%%%%%

%----------------------------------------------------------------------------------------
%	VARIOUS REQUIRED PACKAGES AND CONFIGURATIONS
%----------------------------------------------------------------------------------------

\usepackage[top=3cm,bottom=3cm,left=3cm,right=3cm,headsep=10pt,a4paper]{geometry} % Page margins

\usepackage{graphicx} % Required for including pictures
\graphicspath{{Pictures/}} % Specifies the directory where pictures are stored
\usepackage{tikz} % Required for drawing custom shapes
\usepackage[spanish]{babel} % English language/hyphenation
\usepackage{enumitem} % Customize lists
\setlist{nolistsep} % Reduce spacing between bullet points and numbered lists
\usepackage{booktabs} % Required for nicer horizontal rules in tables
\usepackage{xcolor} % Required for specifying colors by name
\definecolor{ocre}{RGB}{243,102,25} % Define the orange color used for highlighting throughout the book

%----------------------------------------------------------------------------------------
%	FONTS
%----------------------------------------------------------------------------------------

\usepackage{avant} % Use the Avantgarde font for headings
%\usepackage{times} % Use the Times font for headings
\usepackage{mathptmx} % Use the Adobe Times Roman as the default text font together with math symbols from the Sym­bol, Chancery and Com­puter Modern fonts

\usepackage{microtype} % Slightly tweak font spacing for aesthetics
\usepackage[utf8]{inputenc} % Required for including letters with accents
\usepackage[T1]{fontenc} % Use 8-bit encoding that has 256 glyphs

%----------------------------------------------------------------------------------------
%	BIBLIOGRAPHY AND INDEX
%----------------------------------------------------------------------------------------

\usepackage[style=alphabetic,citestyle=numeric,sorting=nyt,sortcites=true,autopunct=true,hyperref=true,abbreviate=false,backref=true,backend=biber]{biblatex}
\addbibresource{bibliography.bib} % BibTeX bibliography file
\defbibheading{bibempty}{}

\usepackage{calc} % For simpler calculation - used for spacing the index letter headings correctly
\usepackage{makeidx} % Required to make an index
\makeindex % Tells LaTeX to create the files required for indexing

%----------------------------------------------------------------------------------------
%	MAIN TABLE OF CONTENTS
%----------------------------------------------------------------------------------------

\usepackage{titletoc} % Required for manipulating the table of contents

\contentsmargin{0cm} % Removes the default margin

% Part text styling
\titlecontents{part}[0cm]
{\addvspace{20pt}\centering\large\bfseries}
{}
{}
{}

% Chapter text styling
\titlecontents{chapter}[1.25cm] % Indentation
{\addvspace{12pt}\large\sffamily\bfseries} % Spacing and font options for chapters
{\color{ocre!60}\contentslabel[\Large\thecontentslabel]{1.25cm}\color{ocre}} % Chapter number
{\color{ocre}}  
{\color{ocre!60}\normalsize\;\titlerule*[.5pc]{.}\;\thecontentspage} % Page number

% Section text styling
\titlecontents{section}[1.25cm] % Indentation
{\addvspace{3pt}\sffamily\bfseries} % Spacing and font options for sections
{\contentslabel[\thecontentslabel]{1.25cm}} % Section number
{}
{\hfill\color{black}\thecontentspage} % Page number
[]

% Subsection text styling
\titlecontents{subsection}[1.25cm] % Indentation
{\addvspace{1pt}\sffamily\small} % Spacing and font options for subsections
{\contentslabel[\thecontentslabel]{1.25cm}} % Subsection number
{}
{\ \titlerule*[.5pc]{.}\;\thecontentspage} % Page number
[]

% List of figures
\titlecontents{figure}[0em]
{\addvspace{-5pt}\sffamily}
{\thecontentslabel\hspace*{1em}}
{}
{\ \titlerule*[.5pc]{.}\;\thecontentspage}
[]

% List of tables
\titlecontents{table}[0em]
{\addvspace{-5pt}\sffamily}
{\thecontentslabel\hspace*{1em}}
{}
{\ \titlerule*[.5pc]{.}\;\thecontentspage}
[]

%----------------------------------------------------------------------------------------
%	MINI TABLE OF CONTENTS IN PART HEADS
%----------------------------------------------------------------------------------------

% Chapter text styling
\titlecontents{lchapter}[0em] % Indenting
{\addvspace{15pt}\large\sffamily\bfseries} % Spacing and font options for chapters
{\color{ocre}\contentslabel[\Large\thecontentslabel]{1.25cm}\color{ocre}} % Chapter number
{}  
{\color{ocre}\normalsize\sffamily\bfseries\;\titlerule*[.5pc]{.}\;\thecontentspage} % Page number

% Section text styling
\titlecontents{lsection}[0em] % Indenting
{\sffamily\small} % Spacing and font options for sections
{\contentslabel[\thecontentslabel]{1.25cm}} % Section number
{}
{}

% Subsection text styling
\titlecontents{lsubsection}[.5em] % Indentation
{\normalfont\footnotesize\sffamily} % Font settings
{}
{}
{}

%----------------------------------------------------------------------------------------
%	PAGE HEADERS
%----------------------------------------------------------------------------------------

\usepackage{fancyhdr} % Required for header and footer configuration

\pagestyle{fancy}
\renewcommand{\chaptermark}[1]{\markboth{\sffamily\normalsize\bfseries\chaptername\ \thechapter.\ #1}{}} % Chapter text font settings
\renewcommand{\sectionmark}[1]{\markright{\sffamily\normalsize\thesection\hspace{5pt}#1}{}} % Section text font settings
\fancyhf{} \fancyhead[LE,RO]{\sffamily\normalsize\thepage} % Font setting for the page number in the header
\fancyhead[LO]{\rightmark} % Print the nearest section name on the left side of odd pages
\fancyhead[RE]{\leftmark} % Print the current chapter name on the right side of even pages
\renewcommand{\headrulewidth}{0.5pt} % Width of the rule under the header
\addtolength{\headheight}{2.5pt} % Increase the spacing around the header slightly
\renewcommand{\footrulewidth}{0pt} % Removes the rule in the footer
\fancypagestyle{plain}{\fancyhead{}\renewcommand{\headrulewidth}{0pt}} % Style for when a plain pagestyle is specified

% Removes the header from odd empty pages at the end of chapters
\makeatletter
\renewcommand{\cleardoublepage}{
\clearpage\ifodd\c@page\else
\hbox{}
\vspace*{\fill}
\thispagestyle{empty}
\newpage
\fi}

%----------------------------------------------------------------------------------------
%	THEOREM STYLES
%----------------------------------------------------------------------------------------

\usepackage{amsmath,amsfonts,amssymb,amsthm} % For math equations, theorems, symbols, etc

\newcommand{\intoo}[2]{\mathopen{]}#1\,;#2\mathclose{[}}
\newcommand{\ud}{\mathop{\mathrm{{}d}}\mathopen{}}
\newcommand{\intff}[2]{\mathopen{[}#1\,;#2\mathclose{]}}
\newtheorem{notation}{Notation}[chapter]

% Boxed/framed environments
\newtheoremstyle{ocrenumbox}% % Theorem style name
{0pt}% Space above
{0pt}% Space below
{\normalfont}% % Body font
{}% Indent amount
{\small\bf\sffamily\color{ocre}}% % Theorem head font
{\;}% Punctuation after theorem head
{0.25em}% Space after theorem head
{\small\sffamily\color{ocre}\thmname{#1}\nobreakspace\thmnumber{\@ifnotempty{#1}{}\@upn{#2}}% Theorem text (e.g. Theorem 2.1)
\thmnote{\nobreakspace\the\thm@notefont\sffamily\bfseries\color{black}---\nobreakspace#3.}} % Optional theorem note
\renewcommand{\qedsymbol}{$\blacksquare$}% Optional qed square

\newtheoremstyle{blacknumex}% Theorem style name
{5pt}% Space above
{5pt}% Space below
{\normalfont}% Body font
{} % Indent amount
{\small\bf\sffamily}% Theorem head font
{\;}% Punctuation after theorem head
{0.25em}% Space after theorem head
{\small\sffamily{\tiny\ensuremath{\blacksquare}}\nobreakspace\thmname{#1}\nobreakspace\thmnumber{\@ifnotempty{#1}{}\@upn{#2}}% Theorem text (e.g. Theorem 2.1)
\thmnote{\nobreakspace\the\thm@notefont\sffamily\bfseries---\nobreakspace#3.}}% Optional theorem note

\newtheoremstyle{blacknumbox} % Theorem style name
{0pt}% Space above
{0pt}% Space below
{\normalfont}% Body font
{}% Indent amount
{\small\bf\sffamily}% Theorem head font
{\;}% Punctuation after theorem head
{0.25em}% Space after theorem head
{\small\sffamily\thmname{#1}\nobreakspace\thmnumber{\@ifnotempty{#1}{}\@upn{#2}}% Theorem text (e.g. Theorem 2.1)
\thmnote{\nobreakspace\the\thm@notefont\sffamily\bfseries---\nobreakspace#3.}}% Optional theorem note

% Non-boxed/non-framed environments
\newtheoremstyle{ocrenum}% % Theorem style name
{5pt}% Space above
{5pt}% Space below
{\normalfont}% % Body font
{}% Indent amount
{\small\bf\sffamily\color{ocre}}% % Theorem head font
{\;}% Punctuation after theorem head
{0.25em}% Space after theorem head
{\small\sffamily\color{ocre}\thmname{#1}\nobreakspace\thmnumber{\@ifnotempty{#1}{}\@upn{#2}}% Theorem text (e.g. Theorem 2.1)
\thmnote{\nobreakspace\the\thm@notefont\sffamily\bfseries\color{black}---\nobreakspace#3.}} % Optional theorem note
\renewcommand{\qedsymbol}{$\blacksquare$}% Optional qed square
\makeatother

% Defines the theorem text style for each type of theorem to one of the three styles above
\newcounter{dummy} 
\numberwithin{dummy}{section}
\theoremstyle{ocrenumbox}
\newtheorem{theoremeT}[dummy]{Teorema}
\newtheorem{problem}{Problem}[chapter]
\newtheorem{exerciseT}{Ejercicio}[chapter]
\theoremstyle{blacknumex}
\newtheorem{exampleT}{Example}[chapter]
\theoremstyle{blacknumbox}
\newtheorem{vocabulary}{Vocabulary}[chapter]
\newtheorem{definitionT}{Definición}[section]
\newtheorem{corollaryT}[dummy]{Corolario}
\theoremstyle{ocrenum}
\newtheorem{proposition}[dummy]{Proposición}

%----------------------------------------------------------------------------------------
%	DEFINITION OF COLORED BOXES
%----------------------------------------------------------------------------------------

\RequirePackage[framemethod=default]{mdframed} % Required for creating the theorem, definition, exercise and corollary boxes

% Theorem box
\newmdenv[skipabove=7pt,
skipbelow=7pt,
backgroundcolor=black!5,
linecolor=ocre,
innerleftmargin=5pt,
innerrightmargin=5pt,
innertopmargin=5pt,
leftmargin=0cm,
rightmargin=0cm,
innerbottommargin=5pt]{tBox}

% Exercise box	  
\newmdenv[skipabove=7pt,
skipbelow=7pt,
rightline=false,
leftline=true,
topline=false,
bottomline=false,
backgroundcolor=ocre!10,
linecolor=ocre,
innerleftmargin=5pt,
innerrightmargin=5pt,
innertopmargin=5pt,
innerbottommargin=5pt,
leftmargin=0cm,
rightmargin=0cm,
linewidth=4pt]{eBox}	

% Definition box
\newmdenv[skipabove=7pt,
skipbelow=7pt,
rightline=false,
leftline=true,
topline=false,
bottomline=false,
linecolor=ocre,
innerleftmargin=5pt,
innerrightmargin=5pt,
innertopmargin=0pt,
leftmargin=0cm,
rightmargin=0cm,
linewidth=4pt,
innerbottommargin=0pt]{dBox}	

% Corollary box
\newmdenv[skipabove=7pt,
skipbelow=7pt,
rightline=false,
leftline=true,
topline=false,
bottomline=false,
linecolor=gray,
backgroundcolor=black!5,
innerleftmargin=5pt,
innerrightmargin=5pt,
innertopmargin=5pt,
leftmargin=0cm,
rightmargin=0cm,
linewidth=4pt,
innerbottommargin=5pt]{cBox}

% Creates an environment for each type of theorem and assigns it a theorem text style from the "Theorem Styles" section above and a colored box from above
\newenvironment{theorem}{\begin{tBox}\begin{theoremeT}}{\end{theoremeT}\end{tBox}}
\newenvironment{exercise}{\begin{eBox}\begin{exerciseT}}{\hfill{\color{ocre}\tiny\ensuremath{\blacksquare}}\end{exerciseT}\end{eBox}}				  
\newenvironment{definition}{\begin{dBox}\begin{definitionT}}{\end{definitionT}\end{dBox}}	
\newenvironment{example}{\begin{exampleT}}{\hfill{\tiny\ensuremath{\blacksquare}}\end{exampleT}}		
\newenvironment{corollary}{\begin{cBox}\begin{corollaryT}}{\end{corollaryT}\end{cBox}}	

%----------------------------------------------------------------------------------------
%	REMARK ENVIRONMENT
%----------------------------------------------------------------------------------------

\newenvironment{remark}{\par\vspace{10pt}\small % Vertical white space above the remark and smaller font size
\begin{list}{}{
\leftmargin=35pt % Indentation on the left
\rightmargin=25pt}\item\ignorespaces % Indentation on the right
\makebox[-2.5pt]{\begin{tikzpicture}[overlay]
\node[draw=ocre!60,line width=1pt,circle,fill=ocre!25,font=\sffamily\bfseries,inner sep=2pt,outer sep=0pt] at (-15pt,0pt){\textcolor{ocre}{R}};\end{tikzpicture}} % Orange R in a circle
\advance\baselineskip -1pt}{\end{list}\vskip5pt} % Tighter line spacing and white space after remark

%----------------------------------------------------------------------------------------
%	SECTION NUMBERING IN THE MARGIN
%----------------------------------------------------------------------------------------

\makeatletter
\renewcommand{\@seccntformat}[1]{\llap{\textcolor{ocre}{\csname the#1\endcsname}\hspace{1em}}}                    
\renewcommand{\section}{\@startsection{section}{1}{\z@}
{-4ex \@plus -1ex \@minus -.4ex}
{1ex \@plus.2ex }
{\normalfont\large\sffamily\bfseries}}
\renewcommand{\subsection}{\@startsection {subsection}{2}{\z@}
{-3ex \@plus -0.1ex \@minus -.4ex}
{0.5ex \@plus.2ex }
{\normalfont\sffamily\bfseries}}
\renewcommand{\subsubsection}{\@startsection {subsubsection}{3}{\z@}
{-2ex \@plus -0.1ex \@minus -.2ex}
{.2ex \@plus.2ex }
{\normalfont\small\sffamily\bfseries}}                        
\renewcommand\paragraph{\@startsection{paragraph}{4}{\z@}
{-2ex \@plus-.2ex \@minus .2ex}
{.1ex}
{\normalfont\small\sffamily\bfseries}}

%----------------------------------------------------------------------------------------
%	PART HEADINGS
%----------------------------------------------------------------------------------------

% numbered part in the table of contents
\newcommand{\@mypartnumtocformat}[2]{%
\setlength\fboxsep{0pt}%
\noindent\colorbox{ocre!20}{\strut\parbox[c][.7cm]{\ecart}{\color{ocre!70}\Large\sffamily\bfseries\centering#1}}\hskip\esp\colorbox{ocre!40}{\strut\parbox[c][.7cm]{\linewidth-\ecart-\esp}{\Large\sffamily\centering#2}}}%
%%%%%%%%%%%%%%%%%%%%%%%%%%%%%%%%%%
% unnumbered part in the table of contents
\newcommand{\@myparttocformat}[1]{%
\setlength\fboxsep{0pt}%
\noindent\colorbox{ocre!40}{\strut\parbox[c][.7cm]{\linewidth}{\Large\sffamily\centering#1}}}%
%%%%%%%%%%%%%%%%%%%%%%%%%%%%%%%%%%
\newlength\esp
\setlength\esp{4pt}
\newlength\ecart
\setlength\ecart{1.2cm-\esp}
\newcommand{\thepartimage}{}%
\newcommand{\partimage}[1]{\renewcommand{\thepartimage}{#1}}%
\def\@part[#1]#2{%
\ifnum \c@secnumdepth >-2\relax%
\refstepcounter{part}%
\addcontentsline{toc}{part}{\texorpdfstring{\protect\@mypartnumtocformat{\thepart}{#1}}{\partname~\thepart\ ---\ #1}}
\else%
\addcontentsline{toc}{part}{\texorpdfstring{\protect\@myparttocformat{#1}}{#1}}%
\fi%
\startcontents%
\markboth{}{}%
{\thispagestyle{empty}%
\begin{tikzpicture}[remember picture,overlay]%
\node at (current page.north west){\begin{tikzpicture}[remember picture,overlay]%	
\fill[ocre!20](0cm,0cm) rectangle (\paperwidth,-\paperheight);
\node[anchor=north] at (4cm,-3.25cm){\color{ocre!40}\fontsize{220}{100}\sffamily\bfseries\@Roman\c@part}; 
\node[anchor=south east] at (\paperwidth-1cm,-\paperheight+1cm){\parbox[t][][t]{8.5cm}{
\printcontents{l}{0}{\setcounter{tocdepth}{1}}%
}};
\node[anchor=north east] at (\paperwidth-1.5cm,-3.25cm){\parbox[t][][t]{15cm}{\strut\raggedleft\color{white}\fontsize{30}{30}\sffamily\bfseries#2}};
\end{tikzpicture}};
\end{tikzpicture}}%
\@endpart}
\def\@spart#1{%
\startcontents%
\phantomsection
{\thispagestyle{empty}%
\begin{tikzpicture}[remember picture,overlay]%
\node at (current page.north west){\begin{tikzpicture}[remember picture,overlay]%	
\fill[ocre!20](0cm,0cm) rectangle (\paperwidth,-\paperheight);
\node[anchor=north east] at (\paperwidth-1.5cm,-3.25cm){\parbox[t][][t]{15cm}{\strut\raggedleft\color{white}\fontsize{30}{30}\sffamily\bfseries#1}};
\end{tikzpicture}};
\end{tikzpicture}}
\addcontentsline{toc}{part}{\texorpdfstring{%
\setlength\fboxsep{0pt}%
\noindent\protect\colorbox{ocre!40}{\strut\protect\parbox[c][.7cm]{\linewidth}{\Large\sffamily\protect\centering #1\quad\mbox{}}}}{#1}}%
\@endpart}
\def\@endpart{\vfil\newpage
\if@twoside
\if@openright
\null
\thispagestyle{empty}%
\newpage
\fi
\fi
\if@tempswa
\twocolumn
\fi}

%----------------------------------------------------------------------------------------
%	CHAPTER HEADINGS
%----------------------------------------------------------------------------------------

\newcommand{\thechapterimage}{}%
\newcommand{\chapterimage}[1]{\renewcommand{\thechapterimage}{#1}}%
\def\@makechapterhead#1{%
{\parindent \z@ \raggedright \normalfont
\ifnum \c@secnumdepth >\m@ne
\if@mainmatter
\begin{tikzpicture}[remember picture,overlay]
\node at (current page.north west)
{\begin{tikzpicture}[remember picture,overlay]
\node[anchor=north west,inner sep=0pt] at (0,0) {\includegraphics[width=\paperwidth]{\thechapterimage}};
\draw[anchor=west] (\Gm@lmargin,-9cm) node [line width=2pt,rounded corners=15pt,draw=ocre,fill=white,fill opacity=0.5,inner sep=15pt]{\strut\makebox[22cm]{}};
\draw[anchor=west] (\Gm@lmargin+.3cm,-9cm) node {\huge\sffamily\bfseries\color{black}\thechapter. #1\strut};
\end{tikzpicture}};
\end{tikzpicture}
\else
\begin{tikzpicture}[remember picture,overlay]
\node at (current page.north west)
{\begin{tikzpicture}[remember picture,overlay]
\node[anchor=north west,inner sep=0pt] at (0,0) {\includegraphics[width=\paperwidth]{\thechapterimage}};
\draw[anchor=west] (\Gm@lmargin,-9cm) node [line width=2pt,rounded corners=15pt,draw=ocre,fill=white,fill opacity=0.5,inner sep=15pt]{\strut\makebox[22cm]{}};
\draw[anchor=west] (\Gm@lmargin+.3cm,-9cm) node {\huge\sffamily\bfseries\color{black}#1\strut};
\end{tikzpicture}};
\end{tikzpicture}
\fi\fi\par\vspace*{270\p@}}}

%-------------------------------------------

\def\@makeschapterhead#1{%
\begin{tikzpicture}[remember picture,overlay]
\node at (current page.north west)
{\begin{tikzpicture}[remember picture,overlay]
\node[anchor=north west,inner sep=0pt] at (0,0) {\includegraphics[width=\paperwidth]{\thechapterimage}};
\draw[anchor=west] (\Gm@lmargin,-9cm) node [line width=2pt,rounded corners=15pt,draw=ocre,fill=white,fill opacity=0.5,inner sep=15pt]{\strut\makebox[22cm]{}};
\draw[anchor=west] (\Gm@lmargin+.3cm,-9cm) node {\huge\sffamily\bfseries\color{black}#1\strut};
\end{tikzpicture}};
\end{tikzpicture}
\par\vspace*{270\p@}}
\makeatother

%----------------------------------------------------------------------------------------
%	HYPERLINKS IN THE DOCUMENTS
%----------------------------------------------------------------------------------------

\usepackage[hyperindex,bookmarks,hidelinks]{hyperref}
\hypersetup{colorlinks=false,breaklinks=true,urlcolor=ocre,bookmarksopen=false,pdftitle={Title},pdfauthor={Author}}
\usepackage{bookmark}
\bookmarksetup{
open,
numbered,
addtohook={%
\ifnum\bookmarkget{level}=0 % chapter
\bookmarksetup{bold}%
\fi
\ifnum\bookmarkget{level}=-1 % part
\bookmarksetup{color=ocre,bold}%
\fi
}
} % Insert the commands.tex file which contains the majority of the structure behind the template

\begin{document}

%----------------------------------------------------------------------------------------
%	TÍTULO
%----------------------------------------------------------------------------------------

\begingroup
\thispagestyle{empty}
\begin{tikzpicture}[remember picture,overlay]
\coordinate [below=12cm] (midpoint) at (current page.north);
\node at (current page.north west)
{\begin{tikzpicture}[remember picture,overlay]
\node[anchor=north west,inner sep=0pt] at (0,0) {\includegraphics[width=\paperwidth]{background}}; % Background image
\draw[anchor=north] (midpoint) node [fill=ocre!30!white,fill opacity=0.6,text opacity=1,inner sep=1cm]{\Huge\centering\bfseries\sffamily\parbox[c][][t]{\paperwidth}
{\centering Teoría de Categorías\\[15pt] % Book title
{\Large Introducción y aplicación a la programación}\\[20pt] % Subtitle
{\huge Mario Román}}}; % Author name
\end{tikzpicture}};
\end{tikzpicture}
\vfill
\endgroup

%----------------------------------------------------------------------------------------
%	COPYRIGHT
%----------------------------------------------------------------------------------------

\newpage
~\vfill
\thispagestyle{empty}

\noindent Copyright \copyright\ 2015 Mario Román\\ % Copyright notice

\noindent \textsc{Doble Grado en Matemáticas e Informática de la Universidad de Granada}\\ % Publisher

\noindent \textsc{dgiim.github.io}\\ % URL

\noindent Licensed under the Creative Commons Attribution-NonCommercial 3.0 Unported License 
(the ``License''). You may not use this file except in compliance with the License. You may 
obtain a copy of the License at \url{http://creativecommons.org/licenses/by-nc/3.0}. Unless 
required by applicable law or agreed to in writing, software distributed under the License 
is distributed on an \textsc{``as is'' basis, without warranties or conditions of any kind}, 
either express or implied. See the License for the specific language governing permissions 
and limitations under the License.\\ % License information

\noindent \textit{First printing, March 2015} % Printing/edition date

%----------------------------------------------------------------------------------------
%	CONTENIDOS
%----------------------------------------------------------------------------------------

\chapterimage{chapter_head_1.pdf} % Table of contents heading image
\pagestyle{empty} % No headers
\tableofcontents % Print the table of contents itself
\cleardoublepage % Forces the first chapter to start on an odd page so it's on the right
\pagestyle{fancy} % Print headers again

%----------------------------------------------------------------------------------------
%	PART
%----------------------------------------------------------------------------------------

\part{Teoría de Categorías}

%----------------------------------------------------------------------------------------
%	CHAPTER 1
%----------------------------------------------------------------------------------------

\chapterimage{chapter_head_2.pdf} % Chapter heading image

\chapter{Categorías}
  \section{Motivación}\index{Motivación}
    Varias estructuras matemáticas (grupos, espacios vectoriales, espacios topológicos \dots) cuentan
    con morfismos que preservan las estructura subyacentes entre ellas. Como ejemplos:
    \begin {center}
    \begin{tabular}{l|l}
      Conjunto & Morfismos \\
      \hline
      Grupos & Homomorfismos de grupos \\
      Espacios topológicos & Funciones continuas \\
      Espacios métricos & Funciones cortas \\
      Conjuntos & Funciones \\
      Espacios vectoriales sobre $\mathbb{K}$ & Funciones lineales sobre $\mathbb{K}$ \\
    \end{tabular}
    \end{center}
    Si estudiamos axiomáticamente las propiedades abstractas de estas estructuras y sus morfismos,
    obtendremos teoremas particularizables a todos estos casos, útiles por sí mismos.
    Una categoría la formarán una clase de estos espacios con estructura y los morfismos entre estos
    espacios; y los teoremas que deduzcamos para todas las categorías podrán aplicarse a cada uno de
    los espacios.

  \section {Definición formal}
    \begin{definition}
    Una \textbf{categoría} $\C$ está definida por:
    \begin{itemize}
      \item Una clase de objetos de la categoría, $Obj(\mathcal{C})$.
      \item Un conjunto de morfismos $Hom_{\C}(A,B)$, poblado o no, entre cada par de objetos $A,B \in Obj(\C)$.
    \end{itemize}
    \end{definition}

    Cumpliendo sus morfismos las siguientes propiedades:
    \begin{itemize}
      \item Para dos morfismos $f \in Hom(A,B)$, $g \in Hom(B,C)$, existe su morfismo composición $f \circ g$.
      \item La composición es asociativa: $ f \circ (g \circ h) = (f \circ g) \circ h$
      \item Todos los objetos tienen un morfismo identidad, $1_{A} \in Hom(A,A)$, 
	  neutro para la composición: $\forall f \in Hom(A,B): f \circ 1_{A} = 1_{B} \circ f = f$
    \end{itemize}
    
    \begin{exercise}
      Demostrar que la identidad es el único elemento neutro para la composición.
    \end{exercise}
  
    \begin{center}
    	
    
    \begin{tabular}{ccc}
        \begin{tikzpicture}[descr/.style={fill=white,inner sep=2.5pt}]
	  \matrix (m) [matrix of math nodes, row sep=3em, column sep=3em]
	  { A & B \\
	     & C \\ };
	  \path[->,font=\scriptsize]
	  (m-1-1) edge node[auto] {$ f $} (m-1-2)
	  (m-1-2) edge node[auto] {$ g $} (m-2-2)
	  (m-1-1) edge node[auto,swap] {$ g \circ f $} (m-2-2);
	\end{tikzpicture} &
%         \begin{tikzcd}[inner sep=2.5pt, row sep=3em, column sep=3em]
%           A \arrow{r}{f} \arrow{rd}[swap]{g \circ f} & B \arrow{d}{g} \\
%           & C
%         \end{tikzcd} &
        
        \begin{tikzpicture}[descr/.style={fill=white,inner sep=2.5pt}]
	  \matrix (m) [matrix of math nodes, row sep=3em, column sep=3em]
	  { A & B & \\ & C & D \\ };
	  \path[->,font=\scriptsize]
	  (m-1-1) edge node[auto] {$ f $} (m-1-2)
	  (m-1-2) edge node[auto] {$ g $} (m-2-2)
	  (m-2-2) edge node[auto] {$ h $} (m-2-3)
	  (m-1-1) edge node[auto,swap] {$ g \circ f $} (m-2-2)
	  (m-1-2) edge node[auto] {$ h \circ g $} (m-2-3);
	\end{tikzpicture} &

        
	
        \begin{tikzpicture}[descr/.style={fill=white,inner sep=2.5pt}]
	  \matrix (m) [matrix of math nodes, row sep=3em, column sep=3em]
	  { A \\ };
	  \path[->,font=\scriptsize]
	  (m-1-1) edge[loop above] node[auto] {$ 1_A $} (m-1-1);
	\end{tikzpicture}
    \end{tabular}
    \end{center}
    \smallskip 
    \textit{Diagramas conmutativos de las propiedades básicas.} \\

  \section{Ejemplos de categorías}
    Como idea simplificadora, podemos pensar que los objetos son conjuntos, y los morfismos, funciones
    entre esos conjuntos; de hecho, el primer ejemplo es la definición de ese caso concreto. 
    Este es un buen modelo intuitivo para trabajar con algunas categorías,
    pero presentaremos ejemplos que rechazan esta intuición.
    \subsubsection{Categoría \texttt{Set}}
      La categoría de los conjuntos con las funciones entre conjuntos como morfismos.
      \begin{gather*}
        Obj(\texttt{Set}) = Clase\ de\ todos\ los\ conjuntos \\
        Hom(A,B)= B^A = \{f \;|\; f: A \rightarrow B \}
      \end{gather*}
      Trivialmente es categoría: las funciones se componen en funciones, la composición es
      asociativa y la identidad funciona como se espera, siendo otra función.
      
    \subsubsection{Categoría \texttt{VectR}}
      La categoría de los espacios vectoriales reales con las funciones lineales entre
      espacios vectoriales reales.
      \begin{gather*}
        Obj(\texttt{VectR}) = Clase\ de\ espacios\ vectoriales\ sobre\ \mathbb{R} \\
        Hom(A,B)= \mathcal{L}_{\mathbb{R}}(A,B)
      \end{gather*}
      \begin{exercise} 
	Observar que \texttt{VectR} es una categoría.
      \end{exercise}

    
    \subsubsection{Categoría \texttt{$(S,\sim)$}}
      Cualquier conjunto $S$ que tenga definida una relación de equivalencia $\sim$ tiene
      definida una categoría asociada en la que los objetos son los elementos del conjunto
      y los morfismos sólo representan casos particulares de la relación de equivalencia.
      \textit{En este ejemplo, los morfismos no son funciones y los objetos no son conjuntos,
      rechazando por primera vez la intuición del primer ejemplo.}
      \begin{gather*}
        Obj((S,\sim)) = S
      \end{gather*}
      Hay un morfismo entre dos elementos si y sólo si están relacionados:
      \begin{align*}
        Hom(a,b)= \twopartdef{(a,b)}{a \sim b}{\emptyset}{a \nsim b}
      \end{align*}
      
\chapter{Tipos de morfismos}
  \section {Isomorfismos}
    \begin{definition}
      Un morfismo $f \in Hom(A,B)$ es \textbf{isomorfismo} si existe un morfismo inverso:
	$g \in Hom(B,A)$ cumpliendo:
      \begin{gather*}
	(g \circ f) = 1_A \qquad (f \circ g) = 1_B
      \end{gather*}
    \end{definition}


    
    \begin{exercise} 
      Probar que el inverso, si existe, es único. Lo notaremos como $f^{-1}$. Observar que si
      existe un inverso por la derecha y un inverso por la izquierda, deben ser iguales.
    \end{exercise}

    
    \subsection{Propiedades de isomorfismos}
      \begin{itemize}
       \item La identidad es isomorfismo: $(1)^{-1} = 1$ 
       \item El inverso de un isomorfismo es isomorfismo: $(f^{-1})^{-1} = f$.
       \item La composición de isomorfismos es isomorfismo: $(g \circ f)^{-1} = f^{-1} \circ g^{-1}$
      \end{itemize}

    Dos objetos $A,B$ son \textbf{isomorfos} si existe un isomorfismo en $Hom(A,B)$, la isomorfía se
    nota como $A \cong B$ y es, por las propiedades anteriores, una relación de equivalencia.
      
    Nótese que, precisamente, los isomorfismos de \texttt{Set} son las biyecciones entre conjuntos.
    Los isomorfismos de \texttt{Top} son los homeomorfismos y los isomorfismos de \texttt{Met} son
    las isometrías.
    
  \section {Epimorfismos y monomorfismos}
    \begin{definition} 
      Un morfismo $f \in Hom(A,B)$ es \textbf{monomorfismo} si para cualquier objeto $C$ y
      para cada par de morfismos $g,h: C \rightarrow A$, el morfismo puede \textit{cancelarse} a la izquierda:
      \begin{gather*}
	f \circ g = f \circ h \quad \Rightarrow \quad g = h
      \end{gather*}
    \end{definition}

    Es decir, si el siguiente diagrama conmuta y $f$ es monomorfismo, entonces $g=h$.
    \begin{center}
     \begin{tikzcd}
       C \arrow[bend left]{r}{g} \arrow[bend right]{r}{h} & A \arrow[tail]{r}{f} & B
     \end{tikzcd}
    \end{center}
    
    \begin{definition} 
      Análogamente, un morfismo $f \in Hom(A,B)$ es \textbf{epimorfismo} si para cualquier objeto $C$ y
      para cada par de morfismos $g,h: B \rightarrow C$, el morfismo puede \textit{cancelarse} a la derecha:
      \begin{gather*}
	g \circ f = h \circ f \quad \Rightarrow \quad g = h
      \end{gather*}
    \end{definition}

    De nuevo, si el siguiente diagrama conmuta y $f$ es epimorfismo, entonces $g=h$.
    %% A - B = C %%
    \begin{center}
     \begin{tikzcd}
       A \arrow[two heads]{r}{f} & B \arrow[bend left]{r}{g} \arrow[bend right]{r}{h} & C
     \end{tikzcd}
    \end{center}
    
  \section {Ejemplos en varias categorías}
    \subsection{Categoría \texttt{Set}}
      En la categoría de los conjuntos, los isomorfismos son las funciones con inversa; lo
      que coincide con la definición de isomorfismos entre conjuntos. Dos conjuntos son
      isomorfos si tienen la misma cardinalidad.
      Los monomorfismos se corresponden con las funciones inyectivas y los epimorfismos con
      las funciones sobreyectivas.
    
    \subsection{Categoría \texttt{$(\mathbb{Z},\leq)$}}
      En esta categoría, los únicos isomorfismos son las identidades. Nótese además
      que todos los morfismos son monomorfismos y epimorfismos, y, aun así, no son
      isomorfismos. Esto sirve de nuevo para no aceptar directamente la intuición
      guiada por la clase \texttt{Set}, donde morfismos que sean monomorfismos y
      epimorfismos deben ser isomorfismos.
      Y la composición se define como:
      \begin{align*}
       (a,b) \circ (b,c) = (a,c)
      \end{align*}
      Probar que es categoría se reduce a notar que la composición de morfismos es morfismo (por
      ser la relación transitiva), que la composición es asociativa y que existe el morfismo identidad
      $(a,a)$, por ser la relación reflexiva.
     
    \subsection{Categoría \texttt{$(S,\leq)$}}
      Cualquier conjunto parcialmente ordenado tiene una categoría asociada. Los objetos
      son sus elementos y los morfismos son casos particulares de la relación de orden.
      \begin{gather*}
        Obj((S,\leq)) = S
      \end{gather*}
      Hay un morfismo entre $a$ y $b$ si y sólo si $a \leq b$:
      \begin{align*}
        Hom(a,b)= \twopartdef{(a,b)}{a \leq b}{\emptyset}{\mbox{no}}
      \end{align*}
      La composición se define como anteriormente: $(a,b) \circ (b,c) = (a,c)$, y es
      trivial volver a probar que se trata de una categoría.
      \medskip
      
      Un posible diagrama conmutativo de la categoría \texttt{$(\mathbb{N},\leq)$} sería:
      \begin{center}
        \begin{tikzpicture}[descr/.style={fill=white,inner sep=2.5pt}]
	  \matrix (m) [matrix of math nodes, row sep=3em, column sep=3em]
	  { 1 & 3 & 5 \\ & 6 & 7 \\ };
	  \path[->,font=\scriptsize]
	  (m-1-1) edge node[auto] {$ (1,3) $} (m-1-2)
	  (m-1-2) edge node[auto] {$ (3,6) $} (m-2-2)
	  (m-2-2) edge node[auto] {$ (6,7) $} (m-2-3)
	  (m-1-1) edge node[auto,swap] {$ (1,6) $} (m-2-2)
	  (m-1-2) edge node[auto] {$ (3,5) $} (m-1-3)
	  (m-1-3) edge node[auto] {$ (5,7) $} (m-2-3);
	\end{tikzpicture}
      \end{center}
 

\chapter{Propiedades universales}
  Varias construcciones en las diversas estructuras que motivaron las categorías pueden parecer
  escogidas arbitrariamente. Sin embargo, la definición de las propiedades universales las dejará
  como las únicas cumpliendo una construcción formal no arbitraria. Además, nos permitirá descubrir
  relaciones más profundas entre las construcciones en las distintas categorías.
  
  \section {Objetos terminales}
    \begin{definition}
      El objeto $I \in Obj(\C)$ se dice \textbf{inicial} cuando a cualquier otro objeto llega
      un único morfismo desde él. Es decir:
      \begin{gather*}
	\forall A \in Obj(\C):\quad \#(Hom(I,A)) = 1
      \end{gather*}
    \end{definition}
 

    \begin{definition}
      Análogamente, el objeto $F \in Obj(\C)$ se dice \textbf{final} cuando desde cualquier otro objeto llega
      un único morfismo hacia él. Es decir:
      \begin{gather*}
	\forall A \in Obj(\C):\quad \#(Hom(A,F)) = 1
      \end{gather*}
    \end{definition}
 
    
    Se llama \textbf{objeto terminal} a un objeto inicial o final y \textbf{objeto cero} a un objeto
    terminal y final. Una categoría no tiene por qué
    tener objetos terminales, y estos no tienen por qué ser únicos. Pero serán \textit{esencialmente únicos},
    es decir, si dos objetos son ambos iniciales o ambos finales en una categoría, serán isomorfos,
    como demostramos a continuación.
    
    \begin{theorem} Los objetos iniciales y los objetos finales de una
    categoría son esencialmente únicos:
      \begin{itemize}
	\item Si $I_1,I_2$ son iniciales, $I_1 \cong I_2$
	\item Si $F_1,F_2$ son finales, $F_1 \cong F_2$
      \end{itemize}
    \end{theorem}

    \begin{proof}
      Para $I_1,I_2$ iniciales, debe cumplirse que existe un sólo $f \in Hom(I_1,I_2)$
      y que existe un sólo $g \in Hom(I_2,I_1)$. Como en $Hom(I_1,I_1)$ y en $Hom(I_2,I_2)$ sólo existe un morfismo,
      que debe ser la identidad, tenemos:
      \begin{gather*}
	g \circ f = 1_{I_1} \quad f \circ g = 1_{I_2}
      \end{gather*}
      Por lo que son isomorfismos y $I_1 \cong I_2$. Análogamente se prueba para objetos finales.
    \end{proof}
    
  \section {Ejemplos}
    \subsection {Objetos terminales en \texttt{Set}}
      En la categoría $\mathtt{Set}$, es objeto inicial el conjunto vacío: para cualquier otro
      conjunto $A$, $Hom(\emptyset, A)$ consta sólo de la inclusión (la función vacía), que 
      podemos notar por $(\emptyset)$. \\
      Todos los conjuntos de un sólo elemento son finales. Sólo hay un morfismo $c_\ast \in Hom(A,\{\ast\})$,
      la función constantemente $\ast$. Obsérvese que entre ellos, todos los conjuntos de un elemento
      son isomorfos como afirma el teorema anterior.
    
    \subsection {Objetos terminales en \texttt{Grp}}
      Entre los grupos, un homomorfismo debe llevar siempre la identidad hacia la identidad
      del otro grupo, y además, no existe el grupo vacío. Esto hace que el objeto terminal
      y final de la categoría sea el grupo que tiene como único elemento la identidad, $\{e\}$.
      Hacia cualquier otro grupo existe un único homomorfismo (la imagen de la identidad debe ser
      la identidad), y desde cualquier otro grupo existe un único homomorfismo (la función 
      constantemente $e$).
    
    \subsection {Objetos terminales en $(\mathbb{N},\leq)$}
      En esta categoría, $0$ es trivialmente objeto inicial y no hay objetos finales.
      Nótese además que en $(\mathbb{Z},\leq)$ no habría ningún objeto terminal. En general,
      en una categoría $(S,\leq)$, los objetos terminales son el máximo y el mínimo, que
      no tienen por qué existir.
    
    \subsection {Proyección al cociente}
      Dado un conjunto $A$, con una relación de equivalencia $\sim$ definimos una categoría
      que tenga por objetos las funciones (!) que preserven la relación de equivalencia.
      Es decir, las funciones de la forma:
      \begin{center}
	\begin{tikzcd}%[scale=0.5]
	    Z \\
	    A \arrow{u}{f}
	\end{tikzcd}
      \end{center}
      Para $Z$ conjunto y $f \in Hom(A,Z)$, cumpliendo $[a]_{\sim} = [b]_{\sim} \Rightarrow f(a) = f(b)$.
      
      Los morfismos entre dos objetos $Hom((Z,f), (W,g))$ de esta categoría serán los diagramas
      conmutativos de la siguiente forma.
      \begin{center}
       \begin{tikzcd}
         Z \arrow{r}{\phi} & W \\
         A \arrow{u}{f} \arrow{ur}[swap]{g}
       \end{tikzcd}
      \end{center}
      Donde, para que se cumpla el diagrama, debería tenerse: $\phi \circ f = g$.
      
      Y la composición entre morfismos es el exterior de la composición de los dos diagramas
      conmutativos, que vuelve a ser diagrama conmutativo:
      \begin{center}
	\begin{tabular}{ccccc}
	  
	  $\left(
	  \begin{tikzcd}
	    & W \arrow{d}{\psi} \\
	    A \arrow{ur}{g} \arrow{r}[swap]{h} & V
	  \end{tikzcd}
	  \right)$
	  &
	  $\circ$
	  &
	  $\left(
	  \begin{tikzcd}
	    Z \arrow{r}{\phi} & W \\
	    A \arrow{u}{f} \arrow{ur}[swap]{g}
	  \end{tikzcd}
	  \right)$
	  &
	  $=$
	  &
	  $\left(
	  \begin{tikzcd}
	    Z \arrow[bend left]{dr}{\psi \circ \phi} \\
	    A \arrow{u}{f} \arrow{r}[swap]{h} & V
	  \end{tikzcd}
	  \right)$
	  
	\end{tabular}
      \end{center}
      
      La construcción de esta categoría puede parecer sorprendente y artificial, pero
      pronto usaremos construcciones similares para obtener propiedades básicas.
      Por ahora, basta notar que en esta categoría, la proyección al cociente $(A/_\sim,\sim)$
      es el objeto inicial:
      \begin{center}
       \begin{tikzcd}
         A/_\sim \arrow{r}{\exists! \tilde f} & Z \\
	 A \arrow{u}{\pi_\sim} \arrow{ur}[swap]{f}
       \end{tikzcd}
      \end{center}

      \begin{proof}
	Definimos $\tilde f ([x]_\sim) = f(x)$. Comprobamos que está bien definida porque
	exigíamos a las funciones que cumplieran $[a]_{\sim} = [b]_{\sim} \Rightarrow f(a) = f(b)$.
	
	El diagrama es conmutativo trivialmente: $\tilde f \circ \pi = f$
      \end{proof} 
 

 \section {Productos y coproductos}
  \subsection{Producto}
    El producto de dos objetos en una categoría se definirá como otro objeto tal
    que una pareja de morfismos que vaya desde un tercer objeto hacia los dos objetos pueda descomponerse
    de manera única a través de su objeto producto.
    En términos de diagramas conmutativos, para $A,B \in Obj(\C)$, su objeto producto,
    $A \times B$ cumple que, para cualquier $Z \in Obj(\C)$ que se descomponga así,
    existe un único morfismo $\phi$ que lo lleva al producto:
    \begin{center}
%     \productCD{ 
%       A,B,Z,A \times B,
%       \pi_A,\pi_B,f,g, \exists! \phi
%     }
    \end{center}
    A ese morfismo $\phi$ que es único lo llamaremos $f \times g$.
    
    \begin{example} 
      En la categoría \texttt{Set}, el producto de dos conjuntos es el producto
      usual: $A \times B = \{(a,b) \; | \; a \in A, \; b \in B \}$. Si tenemos otro conjunto $Z$
      con dos funciones $f,g$ hacia $A$ y hacia $B$, existe una única función hacia $A \times B$
      que hace el diagrama conmutativo, la función:
      \begin{gather*}
      (f \times g)(z) = (f(z),g(z)) \quad \forall z \in Z
      \end{gather*}
      Aquí, las funciones $\pi_A$, $\pi_B$, son las proyecciones usuales desde el producto.
      Es decir, el siguiente diagrama es conmutativo:
      \begin{center}
      \productCD { 
	A,B,Z,A \times B,
	\pi_A,\pi_B,f,g,f \times g
      }
      \end{center}
    \end{example}
    
    \begin{exercise} 
      Demostrar que $f \times g$ es efectivamente la única función que hace que el
      diagrama conmute.
    \end{exercise}

    
    Esta primera definición puede ser reescrita de forma que además se demuestre automáticamente
    la unicidad esencial del objeto producto en el caso de existir. Recurrimos para ello a
    definir una categoría parecida a la que usamos para definir la proyección al cociente.
  
  \subsection{Categorías \texorpdfstring{$\C_{A,B}$}{C A,B}}
    Dados objetos $A,B \in Obj(\C)$ para una categoría $\C$. Definimos la categoría $\C_{A,B}$
    como la categoría teniendo por objetos los diagramas de la forma:
    \begin{center}
      \begin{tikzcd}[column sep=tiny]
	& Z \arrow{dl}[swap]{f_A} \arrow{dr}{f_B} & \\
	A & & B
      \end{tikzcd}
    \end{center}
    Donde $Z$ es otro objeto con dos morfismos $f_A,f_B$ hacia $A$ y $B$.
    
    Y teniendo por morfismos las funciones $\phi$ en diagramas conmutativos de la siguiente forma:
    \begin{center}
      \begin{tikzcd}[column sep=tiny]
        & V \arrow[bend right]{ddl}[swap]{g_A} \arrow[bend left]{ddr}{g_B} \arrow{d}{\phi} & \\
	& Z \arrow{dl}[swap]{f_A} \arrow{dr}{f_B} & \\
	A & & B
      \end{tikzcd}
    \end{center}
    
    Que se componen como podría esperarse:
    \begin{center}
    	\begin{tabular}{ccccc}
	
	  $\left(
	    \begin{tikzcd}[column sep=tiny]
	      & V \arrow[bend right]{ddl}[swap]{g_A} \arrow[bend left]{ddr}{g_B} \arrow{d}{\psi} & \\
	      & W \arrow{dl}[swap]{h_A} \arrow{dr}{h_B} & \\
	      A & & B
	    \end{tikzcd}
	  \right)$
	  &
	  $\circ$
	  &
	  $\left(
	    \begin{tikzcd}[column sep=tiny]
	      & Z \arrow[bend right]{ddl}[swap]{f_A} \arrow[bend left]{ddr}{f_B} \arrow{d}{\phi} & \\
	      & V \arrow{dl}[swap]{g_A} \arrow{dr}{g_B} & \\
	      A & & B
	    \end{tikzcd}
	  \right)$
	  &
	  $=$
	  &
	  $\left(
	    \begin{tikzcd}[column sep=tiny]
	      & Z \arrow[bend right]{ddl}[swap]{f_A} \arrow[bend left]{ddr}{f_B} \arrow{d}[description]{\psi \circ \phi} & \\
	      & W \arrow{dl}[swap]{h_A} \arrow{dr}{h_B} & \\
	      A & & B
	    \end{tikzcd}
	  \right)$
	  
	\end{tabular}
    \end{center}
    
    Pues bien, se define el producto como el objeto final de esta categoría.
    \begin{exercise} 
      Comprobar que coincide con la definición anterior y observar que esto
      prueba su unicidad esencial.
    \end{exercise}
    
  \subsection{Coproducto}
    Análogamente, el coproducto se define como el objeto final de la categoría $\C^{A,B}$,
    que se forma intuitivamente invirtiendo la dirección de las flechas de la categoría $\C_{A,B}$.
    Es decir, tendrá objetos de la forma:
    \begin{center}
      \begin{tikzcd}[column sep=tiny]
	& Z & \\
	A \arrow{ur}{f_A} & & B \arrow{ul}[swap]{f_B}
      \end{tikzcd}
    \end{center}
    Y los morfismos y su composición se definirán análogamente.
    El coproducto de $A$ y $B$, será $A \amalg B$ con dos morfismos $i_A,i_B$ cumpliendo
    el diagrama conmutativo siguiente:
    \begin{center}
    \coproductCD{
      A,B,Z,A \amalg B,
      i_A,i_B,f,g, {(f,g)}     
    }
    \end{center}
    
    \begin{example} 
      En la categoría \texttt{Set}, el coproducto es la unión disjunta de dos conjuntos: 
      $A \amalg B \cong \{(a,1) \; | \; a \in A\} \cup \{(b,2) \; | \; b \in B \}$. Si tenemos otro conjunto $Z$
      con dos funciones $f,g$ desde $A$ y desde $B$, existe una única función desde $A \amalg B$
      que hace el diagrama conmutativo, la función:
      \begin{gather*}
      (f,g)((x,n)) = \twopartdef{f(x)}{n = 1}{g(x)}{n = 2}
      \end{gather*}
    \end{example}


    
  \subsection{Casos particulares}
  \subsubsection{Categoría \texttt{Grp}}
    En la categoría de los grupos, el producto será el producto usual conjuntista con el
    producto de elementos definido componente a componente.
    \begin{align*}
     G_1 \times G_2 = \{(a,b)\; |\; a \in A,\; b \in B\} \\
     (a,b) \ast (a',b') = (a \ast a', b \ast b')
    \end{align*}

    Y el coproducto será el producto libre de grupos, $G \ast H$, formado por las palabras
    reducidas de elementos de ambos grupos. El producto de palabras es su yuxtaposición reducida.
    \begin{align*}
      G_1 \ast G_2 = \{a_1b_1a_2b_2\dots a_kb_k\; |\; a_i \in A,\; b_i \in B\} \\
     (a_1b_1a_2b_2\dots a_kb_k) \ast (a'_1b'_1a'_2b'_2\dots a'_kb'_k) = (a_1b_1\dots a_kb_ka'_1b'_1\dots a'_kb'_k)
    \end{align*}
    
  \subsubsection{Categoría \texttt{($S,\leq$)}}
    En la categoría dada por una relación de orden puede demostarse que serán el producto
    y el coproducto el ínfimo y el supremo, respectivamente.
    \begin{center}
      \begin{tabular}{cc}
	\productCD { 
	  a,b,c,{\inf\{a,b\}},
	  , , , ,
	}
	&
	\coproductCD{
	  a,b,c,{\sup\{a,b\}},
	  , , , ,     
	}
      \end{tabular}
    \end{center}
    
   \subsubsection{Categoría \texttt{($\mathcal{P}(\Omega),\subseteq$)}}
    Esta categoría es similar a la anterior, es un caso particular para subconjuntos
    con la inclusión como relación de orden parcial. Un morfismo de $A$ hacia $B$ indica
    que $B \subseteq A$. La intersección y la unión serán
    aquí el producto y coproducto:
    \begin{center}
      \begin{tabular}{cc}
	\productCD { 
	  A,B,C,{A \cap B},
	  , , , ,
	}
	&
	\coproductCD{
	  A,B,C,{A \cup B},
	  , , , ,     
	}
      \end{tabular}
    \end{center}
  
  \subsubsection{Categoría \texttt{Prop}}
    Sea una categoría que tiene por objetos a las proposiciones lógicas. Existe un morfismo
    desde una proposición a otra si la primera implica la otra, $A \Rightarrow B$.
    En esta categoría, el producto y el coproducto son la conjunción y la disyunción lógicas:
    \begin{center}
      \begin{tabular}{cc}
	\productCD { 
	  A,B,Z,A \wedge B,
	  , , , ,
	}
	&
	\coproductCD{
	  A,B,Z,A \vee B,
	  , , , ,     
	}
      \end{tabular}
    \end{center}
    Aquí el producto y coproducto de morfismos (implicaciones) sería de la forma:
    \begin{align*}
     (Z \Rightarrow A) \times (Z \Rightarrow B) & = (Z \Rightarrow A \wedge B) \\
     ((A \Rightarrow Z), (B \Rightarrow Z)) & = (A \vee B) \Rightarrow Z
    \end{align*}
    Nótese que en esta categoría, el objeto inicial es \texttt{False} y el objeto
    inicial es \texttt{True}.
    
  \subsubsection{Categoría \texttt{Hask}}
    Sea ahora la categoría de tipos de Haskell.\footnote{Aunque, hablando estrictamente, 
    no sea ni siquiera una categoría, es interesante trabajar un poco en ella: \url{http://www.haskell.org/haskellwiki/Hask}} Los objetos son los posibles tipos
    en Haskell y un morfismo entre dos objetos es una función que tome uno de los
    tipos y devuelva el otro.
    Tenemos los diagramas siguientes de producto y coproducto:
        \begin{center}
      \begin{tabular}{cc}
	\productCD {
	  \mathtt{a},\mathtt{b},\mathtt{c},{\quad \mathtt{(a,b)} \quad},
	  \mathtt{fst}, \mathtt{snd}, \mathtt{f}, \mathtt{g}, \mathtt{f><g}
	}
	&
	\coproductCD{
	  a,b,c,\mathtt{Either\ a\ b},
	  \mathtt{Left}, \mathtt{Right}, \mathtt{f}, \mathtt{g}, \mathtt{either\ f\ g}    
	}
      \end{tabular}
    \end{center}
    Hay ciertas similitudes entre las categorías de tipos y proposiciones que
    se quedarán fuera de lo que se expone aquí.
 

\chapter {Functores}
  \section{La categoría opuesta}
  Hasta ahora, pudiera parecer que estamos repitiendo trabajo. Todo lo hemos definido dos veces,
  una vez sobre cada dirección. Composición a izquierda y derecha, monomorfismos y epimorfismos,
  objetos iniciales y finales, productos y coproductos. 
  
  Vamos a definir ahora la categoría $\C^{op}$ como la categoría opuesta de $\C$. Intuitivamente,
  es la categoría que se obtiene invirtiendo el sentido de todos los morfismos. Formalmente,
  esta categoría está formada por:
  \begin{itemize}
   \item $Obj(\C^{op}) = Obj(\C)$
   \item $Hom_{\C^{op}}(A,B) = Hom_{\C}(B,A)$, para cualesquiera $A,B \in Obj(\C^{op})$
  \end{itemize}
  La composición de morfismos es aquí trivial si tenemos en cuenta que hemos dado la vuelta
  a los morfismos:
  \begin{gather*}
   \forall f^{op} \in Hom(B,A), g^{op} \in Hom(C,B):\quad f^{op} \circ g^{op} = (g \circ f)^{op}
  \end{gather*}
  Como podemos observar en estos dos diagramas:
  \begin{center}
    \begin{tabular}{cc}
      \begin{tikzcd}
	A \arrow{r}{f} \arrow[bend right]{rr}[swap]{g\circ f} &
	B \arrow{r}{g} &
	C
      \end{tikzcd}
      &
      \begin{tikzcd}
	A  &
	B \arrow{l}[swap]{f^{op}} &
	C \arrow{l}[swap]{g^{op}} \arrow[bend left]{ll}{f^{op} \circ g^{op}}
      \end{tikzcd}
    \end{tabular}
  \end{center}

  \begin{exercise} 
    Comprobar que la categoría opuesta forma una categoría.
  \end{exercise}
  \begin{exercise} 
    Definir el objeto final y el coproducto en términos del inicial y el producto 
    respectivamente, usando la categoría opuesta.
  \end{exercise}
  \begin{exercise} 
    Probar que un monomorfismo en una categoría es epimorfismo en la opuesta
    y que los isomorfismos se mantienen también en la categoría opuesta.
  \end{exercise}

 
  \section {Functores}
  Una aplicación entre categorías, preservando su estructura, es un functor.
  \begin{definition} Sean $\mathcal{A}, \mathcal{B}$ categorías. Un \textbf{functor} 
  $F: \mathcal{A} \rightarrow \mathcal{B}$ está formado por:
  \begin{itemize}
   \item Una función entre objetos $F: Obj(\mathcal{A}) \rightarrow Obj(\mathcal{B})$
   \item Una función entre morfismos $F: Hom_\mathcal{A}(A,A') \rightarrow Hom_\mathcal{B}(F(A),F(A'))$
  \end{itemize}
  Cumpliendo:
  \begin{itemize}
   \item Respeta la composición: $F(f' \circ f) = F(f') \circ F(f)$
   \item Respeta la identidad: $F(1_A) = 1_{F(A)}$
  \end{itemize}
  \end{definition}

  Es decir, el functor lleva objetos de una categoría a objetos de otra y los morfismos
  entre objetos de una categoría a morfismos de entre las imágenes de los objetos en la
  otra categoría.

  \subsection {Ejemplos}
    \subsubsection{Cardinalidad}
      La cardinalidad, que relaciona el orden parcial de los conjuntos finitos con el
      orden total entre los naturales, puede verse como un functor:
      \begin{align*}
	\#&: \mathtt{(FiniteSet,\subseteq)} \rightarrow \mathtt{(\mathbb{N},\leq)} \\
	\#&(A) = (cardinalidad\ de\ A) \\
	\#&(A \subseteq B) = \#(A) \leq \#(B) 
      \end{align*}
    
    \subsubsection{Endofunctores}
      Un endofunctor es aquel que va de una categoría hacia sí misma. Un ejemplo
      puede ser una función lineal con determinante positivo en $\mathbb{Z}$:
      \begin{align*}
       (\lambda)&: \mathtt{(\mathbb{Z},\leq)} \rightarrow \mathtt{(\mathbb{Z},\leq)} \\
       (\lambda)&(n) = \lambda n \\
       (\lambda)&(p \leq q) = (\lambda p \leq \lambda q)
      \end{align*}

  
  \subsection {Functores de olvido}
    Un ejemplo de functores son los informalmente llamados functores de olvido,
    que \textit{olvidan} propiedades de una categoría al llevarla a otra mediante
    una inclusión.
    
    \begin{example} El functor $U: \mathtt{Grp} \rightarrow \mathtt{Set}$ que olvida la
    estructura de grupo. $U(G)$ es el conjunto de elementos del grupo $G$ y $U(\phi)$
    es la aplicación $\phi$ entre conjuntos, sin verla como un homomorfismo.
    \end{example}
    
  \subsection {Functores libres}
    Los functores libres hacen un papel inverso a los functores de olvido.
    Informalmente, llevan una categoría hacia otra con más información dotándola del mínimo
    contexto posible.
    
    \begin{example} El functor $F: \mathtt{Set} \rightarrow \mathtt{Grp}$ lleva cada conjunto $X$
    al grupo libre generado por los elementos de $X$, es decir, al grupo de palabras
    reducidas formadas con elementos de $X$. Las funciones entre conjuntos se convierten
    trivialmente en homomorfismos entre sus grupos libres generados.
    \end{example}

\section {Transformaciones naturales}
  \begin{definition} 
    Una transformación natural $\alpha$ entre dos functores $F,G: \C \rightarrow \mathcal{D}$
    viene determinada por:
    \begin{itemize}
    \item Para cada $X \in \C$, un morfismo: \, $\alpha_X : F(X) \rightarrow G(X)$
    \end{itemize}
    Cumpliendo que:
    \begin{itemize}
    \item Para cualquier morfismo $f \in Hom(X,Y)$ se tenga: \, $\alpha_Y \circ Ff = Gf \circ \alpha_X$
    \end{itemize}
  \end{definition}
  Lo que queda representado en el siguiente diagrama conmutativo de naturalidad:
  \begin{center}
  \begin{tikzpicture}[descr/.style={fill=white,inner sep=2.5pt}]
    \matrix (m) [matrix of math nodes, row sep=3em, column sep=3em]
    { FX & GX \\
      FY & GY \\ };
    \path[->,font=\scriptsize]
    (m-1-1) edge node[auto] {$ \alpha_X $} (m-1-2)
    (m-2-1) edge node[auto] {$ \alpha_Y $} (m-2-2)
    (m-1-2) edge node[auto] {$ Gf $} (m-2-2)
    (m-1-1) edge node[auto,swap] {$ Ff $} (m-2-1);
  \end{tikzpicture}
  \end{center}
 
%----------------------------------------------------------------------------------------
%	CHAPTER 2
%----------------------------------------------------------------------------------------

\chapter{In-text Elements}

\section{Theorems}\index{Theorems}

This is an example of theorems.

\subsection{Several equations}\index{Theorems!Several Equations}
This is a theorem consisting of several equations.

\begin{theorem}[Name of the theorem]
In $E=\mathbb{R}^n$ all norms are equivalent. It has the properties:
\begin{align}
& \big| ||\mathbf{x}|| - ||\mathbf{y}|| \big|\leq || \mathbf{x}- \mathbf{y}||\\
&  ||\sum_{i=1}^n\mathbf{x}_i||\leq \sum_{i=1}^n||\mathbf{x}_i||\quad\text{where $n$ is a finite integer}
\end{align}
\end{theorem}

\subsection{Single Line}\index{Theorems!Single Line}
This is a theorem consisting of just one line.

\begin{theorem}
A set $\mathcal{D}(G)$ in dense in $L^2(G)$, $|\cdot|_0$. 
\end{theorem}

%------------------------------------------------

\section{Definitions}\index{Definitions}

This is an example of a definition. A definition could be mathematical or it could define a concept.

\begin{definition}[Definition name]
Given a vector space $E$, a norm on $E$ is an application, denoted $||\cdot||$, $E$ in $\mathbb{R}^+=[0,+\infty[$ such that:
\begin{align}
& ||\mathbf{x}||=0\ \Rightarrow\ \mathbf{x}=\mathbf{0}\\
& ||\lambda \mathbf{x}||=|\lambda|\cdot ||\mathbf{x}||\\
& ||\mathbf{x}+\mathbf{y}||\leq ||\mathbf{x}||+||\mathbf{y}||
\end{align}
\end{definition}

%------------------------------------------------

\section{Notations}\index{Notations}

\begin{notation}
Given an open subset $G$ of $\mathbb{R}^n$, the set of functions $\varphi$ are:
\begin{enumerate}
\item Bounded support $G$;
\item Infinitely differentiable;
\end{enumerate}
a vector space is denoted by $\mathcal{D}(G)$. 
\end{notation}

%------------------------------------------------

\section{Remarks}\index{Remarks}

This is an example of a remark.

\begin{remark}
The concepts presented here are now in conventional employment in mathematics. Vector spaces are taken over the field $\mathbb{K}=\mathbb{R}$, however, established properties are easily extended to $\mathbb{K}=\mathbb{C}$.
\end{remark}

%------------------------------------------------

\section{Corollaries}\index{Corollaries}

This is an example of a corollary.

\begin{corollary}[Corollary name]
The concepts presented here are now in conventional employment in mathematics. Vector spaces are taken over the field $\mathbb{K}=\mathbb{R}$, however, established properties are easily extended to $\mathbb{K}=\mathbb{C}$.
\end{corollary}

%------------------------------------------------

\section{Propositions}\index{Propositions}

This is an example of propositions.

\subsection{Several equations}\index{Propositions!Several Equations}

\begin{proposition}[Proposition name]
It has the properties:
\begin{align}
& \big| ||\mathbf{x}|| - ||\mathbf{y}|| \big|\leq || \mathbf{x}- \mathbf{y}||\\
&  ||\sum_{i=1}^n\mathbf{x}_i||\leq \sum_{i=1}^n||\mathbf{x}_i||\quad\text{where $n$ is a finite integer}
\end{align}
\end{proposition}

\subsection{Single Line}\index{Propositions!Single Line}

\begin{proposition} 
Let $f,g\in L^2(G)$; if $\forall \varphi\in\mathcal{D}(G)$, $(f,\varphi)_0=(g,\varphi)_0$ then $f = g$. 
\end{proposition}

%------------------------------------------------

\section{Examples}\index{Examples}

This is an example of examples.

\subsection{Equation and Text}\index{Examples!Equation and Text}

\begin{example}
Let $G=\{x\in\mathbb{R}^2:|x|<3\}$ and denoted by: $x^0=(1,1)$; consider the function:
\begin{equation}
f(x)=\left\{\begin{aligned} & \mathrm{e}^{|x|} & & \text{si $|x-x^0|\leq 1/2$}\\
& 0 & & \text{si $|x-x^0|> 1/2$}\end{aligned}\right.
\end{equation}
The function $f$ has bounded support, we can take $A=\{x\in\mathbb{R}^2:|x-x^0|\leq 1/2+\epsilon\}$ for all $\epsilon\in\intoo{0}{5/2-\sqrt{2}}$.
\end{example}

\subsection{Paragraph of Text}\index{Examples!Paragraph of Text}

\begin{example}[Example name]
 Lorem ipsum.
\end{example}

%------------------------------------------------

\section{Exercises}\index{Exercises}

This is an example of an exercise.

\begin{exercise}
This is a good place to ask a question to test learning progress or further cement ideas into students' minds.
\end{exercise}

%------------------------------------------------

\section{Problems}\index{Problems}

\begin{problem}
What is the average airspeed velocity of an unladen swallow?
\end{problem}

%------------------------------------------------

\section{Vocabulary}\index{Vocabulary}

Define a word to improve a students' vocabulary.

\begin{vocabulary}[Word]
Definition of word.
\end{vocabulary}

%----------------------------------------------------------------------------------------
%	PART
%----------------------------------------------------------------------------------------

\part{Part Two}

%----------------------------------------------------------------------------------------
%	CHAPTER 3
%----------------------------------------------------------------------------------------

\chapterimage{chapter_head_1.pdf} % Chapter heading image

\chapter{Presenting Information}

\section{Table}\index{Table}

\centering
\begin{tabular}{l l l}
\toprule
\textbf{Treatments} & \textbf{Response 1} & \textbf{Response 2}\\
\midrule
Treatment 1 & 0.0003262 & 0.562 \\
Treatment 2 & 0.0015681 & 0.910 \\
Treatment 3 & 0.0009271 & 0.296 \\
\bottomrule
\end{tabular}

%------------------------------------------------

\section{Figure}\index{Figure}

\centering\includegraphics[scale=0.5]{placeholder}


%----------------------------------------------------------------------------------------
%	BIBLIOGRAPHY
%----------------------------------------------------------------------------------------

\chapter*{Bibliography}
\addcontentsline{toc}{chapter}{\textcolor{ocre}{Bibliography}}
\section*{Books}
\addcontentsline{toc}{section}{Books}
\printbibliography[heading=bibempty,type=book]
\section*{Articles}
\addcontentsline{toc}{section}{Articles}
\printbibliography[heading=bibempty,type=article]

%----------------------------------------------------------------------------------------
%	INDEX
%----------------------------------------------------------------------------------------

\cleardoublepage
\phantomsection
\setlength{\columnsep}{0.75cm}
\addcontentsline{toc}{chapter}{\textcolor{ocre}{Index}}
\printindex

%----------------------------------------------------------------------------------------

\end{document}